\documentclass[11pt]{article}

\usepackage[margin=1in]{geometry}
\usepackage{amsmath,amssymb}
\usepackage{booktabs}
\usepackage{array}
\usepackage{xcolor}
\usepackage[most]{tcolorbox}
\usepackage{enumitem}
\usepackage[colorlinks=true,linkcolor=blue!50!black,urlcolor=blue!50!black,citecolor=blue!50!black]{hyperref}

\newtcolorbox{takeaway}[1][]{colback=yellow!12,colframe=orange!70!black,
  title=#1,fonttitle=\bfseries,boxrule=0.6pt,arc=2pt,left=4pt,right=4pt,top=3pt,bottom=3pt}

\newcommand{\kl}{\ensuremath{\kappa_\lambda}}
\newcommand{\mhh}{\ensuremath{m_{HH}}}
\newcommand{\bbtt}{\ensuremath{b\bar b\tau\tau}}
\setlist{itemsep=2pt,topsep=3pt}

\title{\vspace{-2.5em}Foundation Models for Precision Classification in NSBI:\\
toward a Higgs self-coupling (\kl) measurement in $HH\to\bbtt$}
\author{\normalsize\itshape working note --- draft for discussion}
\date{19 June 2026}

\begin{document}
\maketitle
\vspace{-1.5em}

\begin{takeaway}[Bottom line]
For constraining \kl, the statistically-limited piece is \emph{kinematic shape}
discrimination of nearly-identical distributions. Both our own tests and the
literature indicate this piece is \textbf{method-independent and statistics-limited}:
a foundation model (FM) will \emph{not} out-classify a good high-level-feature NSBI
on the \kl\ shape \emph{(though a systematics-robust FM can still help the \kl\
inference under profiling --- see \S\ref{sec:axes})}. The genuine leverage of an FM is
elsewhere ---
\textbf{(i) data efficiency} at low signal stats, \textbf{(ii) signal-vs-background
acceptance/purity}, and \textbf{(iii) robustness to systematics}. The recommended
path is a \emph{factorized, calibration-first} strategy: keep the validated
signal-extraction machinery and do NSBI on high-level kinematics now; pursue a
jet-free / event-level FM on the acceptance axis as a staged, calibration-gated
upgrade.
\end{takeaway}

\section{Physics goal and current status}
We aim to constrain the Higgs trilinear self-coupling modifier
$\kl=\lambda_{HHH}/\lambda_{HHH}^{\mathrm{SM}}$ through non-resonant di-Higgs
production in the $HH\to\bbtt$ final state, with the dominant sensitivity in the
\textbf{resolved regime} ($\mhh\in[250,400]\,$GeV). \kl\ enters as a shape
distortion of \mhh\ (interference between the triangle, $\propto\lambda_{HHH}$, and
box amplitudes in ggF).

A neutral-simulation-based-inference (NSBI) fit using \textbf{10 event-level features},
\[
\mhh,\ \cos\theta^\star,\ p_T^{HH},\ m_{bb},\ \Delta R_{bb},\ m_{\tau\tau},\
\Delta\phi_{HH},\ \Delta R_{\tau\tau},\ p_T^{H_1},\ p_T^{H_2},
\]
already \textbf{tightens the \kl\ 95\% CI by roughly a factor of two} (with $t\bar t$
background only, so far). The central \emph{challenge} is statistical: very low
signal yield, made worse once systematic uncertainties are profiled.

\textbf{Hypothesis under test:} a foundation-model + fine-tuning paradigm for the
density-ratio (precision-classification) tasks inside the NSBI could improve the
constraint --- primarily via data efficiency and better object use.

\section{The key decomposition: a \kl\ NSBI is two ML jobs}
The single most clarifying point is that the measurement contains two distinct
learning problems, which an FM helps very differently:

\begin{table}[h]
\centering\small
\begin{tabular}{@{}p{0.20\textwidth}p{0.30\textwidth}p{0.42\textwidth}@{}}
\toprule
\textbf{Job} & \textbf{Discriminates} & \textbf{Information needed / best tool} \\
\midrule
(A) Signal vs.\ background & $HH\to\bbtt$ vs.\ $t\bar t$/QCD/fakes &
\emph{Object identity}: b-tagging, real-vs-fake $\tau$, substructure ---
\textbf{flavor FM (Sophon-ak4) territory}. \\
\addlinespace
(B) \kl\ shape / inference & $\kl=0$ vs.\ $5$ vs.\ $\dots$ (morphing) &
\emph{Event kinematics} (\mhh\ shape) --- high-level features / NSBI; a
``precision classification of nearly-identical distributions'' problem. \\
\bottomrule
\end{tabular}
\end{table}

\noindent The \kl\ \emph{constraint} comes from job (B). Job (A) sets the effective
signal purity and is where the analysis is acceptance- and systematics-limited.
A flavor-pretrained FM is built for (A), not (B) --- which explains the results below.

\section{What we tested, and what we found}
We adapted the \texttt{nsbi-lhc-toolkit} with a vendored Particle Transformer
(ParT) backbone and \textbf{Sophon-ak4} (the AK4/resolved-jet ParT foundation model,
pretrained for jet-flavor tagging on JetClass-II). A hierarchical model encodes each
jet's constituents and aggregates jets+objects with an event set-transformer.
Three controls were compared on the \textbf{$\kl{=}0$ vs.\ $\kl{=}5$} task
(with $\kl{=}5$ as the fixed reference/denominator):
\texttt{scratch} (random-init ParT), \texttt{sophon\_frozen}, \texttt{sophon\_finetune}.

\begin{table}[h]
\centering\small
\caption*{\textbf{Validation weighted-BCE (lower = better; random floor $=\ln 2 = 0.693$).}}
\begin{tabular}{@{}lccc@{}}
\toprule
$N$ / point & \texttt{scratch} & \texttt{sophon\_finetune} & \texttt{sophon\_frozen} \\
\midrule
100k (40 ep)$^\dagger$ & 0.6378 & \textbf{0.6351} & 0.6486 \\
2{,}000          & \textbf{0.6752} & 0.6775 & 0.6798 \\
5{,}000          & \textbf{0.6387} & 0.6457 & 0.6538 \\
\bottomrule
\end{tabular}
\end{table}
\vspace{-0.4em}
{\footnotesize $^\dagger$earlier configuration (before early-stopping and the lowered
fine-tune LR); single seed at low $N$ (error bars pending).}

\begin{takeaway}[Empirical finding]
On the \kl\ \emph{shape} task, pretraining does \textbf{not} help --- the ordering is
consistently \texttt{scratch} $\le$ \texttt{finetune} $<$ \texttt{frozen}. Frozen
flavor features are the worst; fine-tuning merely relearns toward scratch.
\end{takeaway}

\noindent\textbf{Interpretation.} (i)~\emph{Pretext mismatch}: Sophon-ak4 is optimized
for jet flavor, whose per-jet embeddings are lossy about the $p_T$/mass needed to
reconstruct \mhh\ from cross-jet correlations. (ii)~The task is precisely
``precision classification of nearly-identical distributions,'' which (next section)
is statistics-limited and largely architecture-independent.

\section{External evidence and the FM landscape}
A 2025 review~\cite{fmreview} organizes HEP foundation models into \emph{supervised}
(OmniLearn, Sophon, RS3L) and \emph{self-supervised} (OmniJet-$\alpha$, MPM) families.
Points most relevant here:

\begin{itemize}
\item \textbf{Explicit vs.\ implicit priors are a wash --- and statistics-limited}
  \cite{wamorkar}. A direct comparison of L-GATr (Lorentz-equivariant) vs.\ OmniLearn
  (FM pretrained on 100M events) on \emph{precision classification of similar classes}
  finds ``comparable performance \dots given the statistical precision of the
  datasets,'' and that the gains from encoding physics are ``method-independent.''
  This benchmark \emph{is} our \kl-shape task, and predicts our null result. Nuances:
  on the H1/$ep$ task the FM \emph{beat} the equivariant net, and for small signal
  injections the implicit (FM) approach was slightly better --- a mild lean toward
  FMs in the low-statistics corner.
\item \textbf{Taus are not a blocker.} OmniJet-$\alpha$ fine-tuned for hadronic-$\tau$
  reconstruction improved momentum resolution by $\sim$50\% vs.\ from-scratch
  training~\cite{omnijettau} --- so an earlier worry that OmniLearn-style FMs are
  unsuited to $\tau$ final states was premature.
\item \textbf{Event-level FMs exist.} PECM~\cite{pecm} is a GNN over event objects
  (jets, $e$, $\mu$, $\gamma$, MET with four-momenta + type IDs) pretrained on 120M
  events; it reports $>$5 percentage-point accuracy gains at $10^3$--$10^5$
  events/class (our regime) and $\sim$12$\times$ faster fine-tuning, transferring
  across simulators. Closest match to our event-level inputs.
\item \textbf{Self-supervision retains kinematics; RS3L targets systematics.}
  MPM~\cite{mpm} learns task-agnostic representations (less biased than supervised
  flavor labels). RS3L~\cite{rs3l} pretrains on \emph{re-simulation} augmentations
  (varied shower/hadronization/detector), yielding embeddings \emph{robust to those
  systematics} --- directly relevant to a systematics-limited measurement.
\end{itemize}

\section{Two architectural strategies for the measurement}
\textbf{Strategy A --- jet-free, end-to-end.} Feed \emph{all} particle-flow
candidates of the event (no jet clustering) into a ParT/event-FM, add the 10
high-level features, and do NSBI for \kl\ in one model. Motivated by a recent
\emph{jet-free} HH$\to 4b$ analysis~\cite{jetfree4b}: ingesting all PF candidates
(\,$\sim$300/event) into a ParT gave $>$5$\times$ search significance and an
order-of-magnitude more retained signal at high background rejection. Caveats: it is
\emph{not} an FM (task-specific, trained on $\sim$$10^8$ events --- ``training
statistics play a critical role''), it is HH$\to 4b$ (a worse case for jet methods
than \bbtt), and it required a \emph{new} ``calibrate--validate--measure'' scheme
(calibrate on $ZZ\to4b$, validate on $ZH\to4b$) because object-level scale factors do
not apply to a jet-free discriminant.

\textbf{Strategy B --- factorized.} Use the current CMS-style signal extractor for
job (A), then NSBI on the high-level features (or an FM embedding) for job (B). The
CMS $HH\to\bbtt$ analysis~\cite{cmsbbtt} already uses a sophisticated DNN: a Lorentz
Boost Network front-end, a $\tau\tau$ neutrino-regression sub-net, and a DenseNet
classifier, with inputs that are full-system four-momenta \emph{plus}
ParticleNet/UParT (b-tag, c-vs-b, c-vs-l) and HHbTag scores --- i.e.\ it already uses
ParT as fixed tagger \emph{scores} on reconstructed objects, with the entire
validated systematics machinery built around those objects.

\begin{table}[h]
\centering\small
\begin{tabular}{@{}p{0.21\textwidth}p{0.36\textwidth}p{0.36\textwidth}@{}}
\toprule
 & \textbf{A: jet-free FM ($+$10 feats) $\to$ NSBI} & \textbf{B: CMS S/B $\to$ NSBI on feats/FM} \\
\midrule
Information ceiling & Highest (soft radiation, no clustering/WP losses) & High, capped by reconstruction \\
Where it wins & \textbf{Signal acceptance \& purity} & Clean \kl-kinematics (the proven $\sim$2$\times$) \\
\kl-shape info & Same (kinematic) & Same (kinematic) \\
Calibration / systematics & \textbf{Hard \& novel} (new scheme) & \textbf{Validated CMS machinery} \\
NSBI tractability & High-dim, harder closure/morphing & Low-dim, interpretable, easy closure \\
Data efficiency & Needs huge stats unless pretrained & Already efficient \\
Risk / time-to-result & High / frontier & Low / near-term \\
\bottomrule
\end{tabular}
\end{table}

\noindent\textbf{Important caveat:} ``Sophon-ak4 jet-free'' mixes scales ---
Sophon-ak4 was pretrained on \emph{single AK4 jets} ($\sim$128 constituents), so
applying it to the \emph{whole event} ($\sim$300 particles) is out-of-distribution.
Jet-free-whole-event calls for an \emph{event-scale} FM (PECM-style, or a from-scratch
/ SSL event-ParT), not Sophon-ak4 stretched beyond its domain.

\section{Where the FM actually helps: separation vs.\ systematics robustness}
\label{sec:axes}
A natural but \emph{over-simplified} reading of the two-job decomposition is ``the FM
belongs in the $s/b$ ratio, not the \kl\ ratio.'' That holds for one benefit axis and
fails for the other --- and the one it fails for is precisely the systematics-limited
regime that constrains this measurement. An NSBI \kl\ fit needs (at least) two learned
ratios: a signal-vs-background ratio $s(x)/b(x)$ and a \kl-parametrized signal-shape
ratio $r(x;\kl)$. An FM is a \emph{backbone for a ratio estimator}, so the question is
always ``which ratio, and helping \emph{what}?''

\begin{itemize}
\item \textbf{Axis 1 --- separation / central sensitivity.} How well a ratio
  discriminates (the kl0-vs-kl5 ceiling above). Here the 10 features are at ceiling and
  data-efficient, so an FM backbone does \emph{not} improve the \kl\ ratio's
  discrimination; its separation gain accrues only to $s/b$ (acceptance, object ID).
\item \textbf{Axis 2 --- robustness $+$ estimator stability under limited stats.} A
  property of the learned \emph{representation}, not of its separation power. It can
  improve \emph{either} ratio --- including \kl\ --- and it is the axis that matters
  when, as here, the measurement is systematics- and statistics-limited.
\end{itemize}

\begin{table}[h]
\centering\small
\begin{tabular}{@{}lll@{}}
\toprule
 & \textbf{Axis 1: separation} & \textbf{Axis 2: robustness / stats stability} \\
\midrule
$s/b$ ratio & FM helps (acceptance, object ID) & FM helps (robust $+$ data-efficient) \\
\kl\ shape ratio & FM does \emph{not} help (at ceiling) & FM \emph{can} help (robust rep
$\to$ smaller shape syst.) \\
\bottomrule
\end{tabular}
\end{table}

\noindent\textbf{How an FM helps the \kl\ ratio on Axis 2:}
\begin{enumerate}[leftmargin=1.4em]
\item \emph{Robustness} (e.g.\ RS3L~\cite{rs3l}): a representation invariant to
  shower/hadronization/detector variations yields a discriminant whose \emph{shape}
  shifts less under nuisances $\to$ smaller profiled systematics $\to$ tighter CI
  \emph{even at fixed separation}.
\item \emph{Data-efficient systematic templates}: building systematic up/down variations
  needs ratios estimated on stats-starved varied MC; a pretrained, data-efficient
  estimator gives less noisy templates $\to$ better-behaved profiling. (Largest where
  the estimator is high-dimensional; smaller for the already-efficient 10-feature ratio.)
\item \emph{(Advanced) nuisance-parametrized NSBI}: learn $r(x;\kl,\theta_{\rm syst})$
  directly; an FM backbone makes the higher-dimensional conditional learning more
  data-efficient.
\end{enumerate}

\noindent\textbf{Catches (not free):}
\begin{itemize}
\item To get the Axis-2 benefit on \kl\ you build the ratio on the FM representation
  rather than the at-ceiling hand features --- trading interpretability / easy closure
  for robustness.
\item You swap \emph{well-calibrated object} systematics (JES, $\tau$ energy scale) for
  the FM representation's \emph{own} systematics, which must then be established (the
  calibrate--validate--measure burden); robustness to known systematics can introduce
  new ones.
\item Whether it is a \emph{net} win is empirical --- and it \textbf{will not appear in
  BCE/AUC}, which measure separation. We have so far measured only separation, which by
  construction cannot reveal a robustness gain.
\end{itemize}

\begin{takeaway}[Corrected conclusion]
The FM's \emph{separation} benefit accrues only to $s/b$; its \emph{robustness /
data-efficiency} benefit can accrue to \emph{either} ratio, including \kl. So an FM will
not beat the high-level features on \kl\ \emph{separation}, but a systematics-robust /
data-efficient FM (RS3L-style) can still help the \kl\ \emph{measurement} by making the
inference more robust under profiling when stats are limited --- the one place an FM
earns its keep on the kinematic side.
\end{takeaway}

\paragraph{The diagnostic this implies.} The right test of the FM's value here is
\emph{not} classifier loss but \emph{robustness}: propagate the dominant nuisances
(JES, $\tau$ energy scale, parton shower) through each control and measure (i) how much
the discriminant distribution shifts (nominal vs.\ varied) and (ii) whether the
\emph{profiled} \kl\ CI shrinks. A feature-based vs.\ FM-based comparison on \emph{those}
metrics --- not BCE/AUC --- is what would show an FM helping where the analysis is
actually limited.

\section{Synthesis and recommendation}
\begin{takeaway}[Key insights]
\begin{enumerate}[leftmargin=1.4em]
\item The \kl\ constraint is \textbf{kinematic and statistics-limited}; explicit
  (equivariant) and implicit (FM) priors give comparable shape performance. Do not
  expect an FM to beat the high-level-feature NSBI on the \kl\ shape.
\item The FM's real value is \textbf{data efficiency}, \textbf{S/B acceptance/purity},
  and \textbf{systematics robustness} --- not shape extraction.
\item For a systematics-limited result, \textbf{calibration of any learned
  discriminant is the dominant risk}; the jet-free upside can be eaten by a poorly
  understood systematic.
\item Use the right FM for the right job: Sophon-ak4 (jet-scale) $\to$ per-jet S/B;
  event-scale FM $\to$ jet-free acceptance; high-level features / NSBI $\to$ \kl.
\end{enumerate}
\end{takeaway}

\noindent\textbf{Recommended staged plan.}
\begin{itemize}
\item \textbf{Near-term (low risk):} Strategy B, sharpened. Do NSBI on the high-level
  features (optionally $+$ the CMS DNN score) within the validated signal regions ---
  the direct productization of the $\sim$2$\times$ result on real, calibrated data.
\item \textbf{Test the right task:} a signal-vs-$t\bar t$ diagnostic (3 controls) ---
  where per-jet Sophon \emph{should} help, unlike the \kl-shape task. Plus a
  high-level kinematic ceiling baseline (BDT) to bound how much shape headroom exists.
\item \textbf{Parallel R\&D (high ceiling):} a jet-free / \emph{event-level} FM for
  acceptance, introduced as a \emph{calibratable} added discriminant (not a
  replacement), folded into the NSBI only once its calibration is validated.
\item \textbf{Distinctive contribution:} RS3L-style re-simulation pretraining for a
  \emph{systematics-robust} embedding --- aimed squarely at the systematics bottleneck,
  and evaluated via the robustness diagnostic of \S\ref{sec:axes} (nuisance shape
  variation + profiled CI), not BCE/AUC.
\end{itemize}

\section{Open questions}
\begin{itemize}
\item Do explicit and implicit priors learn the \emph{same} physics, and does
  combining them (a Lorentz-equivariant \emph{foundation} model) help in our regime?
\item How large is the jet-free acceptance gain for \bbtt\ (cleaner than $4b$), and
  can it be calibrated with $Z\to\tau\tau$/$Z\to b\bar b$ control samples?
\item Does any S/B or acceptance gain actually translate into a tighter \kl\ CI once
  systematics are profiled (the figure of merit is the constraint, not BCE/AUC)?
\end{itemize}

\begin{thebibliography}{9}\small
\bibitem{fmreview} \emph{Foundation models for high-energy physics} (review),
  arXiv:2509.21434.
\bibitem{wamorkar} T.~Wamorkar \emph{et al.}, \emph{Explicit or Implicit? Encoding
  Physics at the Precision Frontier}, arXiv:2603.08802.
\bibitem{omnijettau} \emph{Reconstructing hadronically decaying tau leptons with a
  jet foundation model} (OmniJet-$\alpha$), arXiv:2503.19165, SciPost Phys.\ Core (2025).
\bibitem{pecm} \emph{Pretrained Event Classification Model for High Energy Physics
  Analysis}, arXiv:2412.10665.
\bibitem{mpm} \emph{Masked Particle Modeling on Sets: Towards Self-Supervised HEP
  Foundation Models}, arXiv:2401.13537.
\bibitem{rs3l} \emph{Re-Simulation-based Self-Supervised Learning for Pre-Training
  Foundation Models}, arXiv:2403.07066; Phys.\ Rev.\ D \textbf{111}, 032010 (2025).
\bibitem{jetfree4b} \emph{A calibratable, jet-free framework for $HH\to4b$ via
  all-particle multiclass classification}, arXiv:2508.15048.
\bibitem{cmsbbtt} CMS Collaboration, \emph{Search for non-resonant Higgs boson pair
  production in the $\bbtt$ final state with Run-3 data}, CMS AN-25-103 (draft, 2026).
\bibitem{omnilearn} O.~Amram, K.~Mikuni, B.~Nachman \emph{et al.}, \emph{OmniLearn};
  \emph{OmniJet-$\alpha$}, arXiv:2403.05618.
\bibitem{lgatr} \emph{A Lorentz-Equivariant Transformer for All of the LHC} (L-GATr),
  arXiv:2411.00446; \emph{Lorentz-Equivariant Geometric Algebra Transformers},
  arXiv:2405.14806.
\end{thebibliography}

\end{document}
